% !TeX spellcheck = es_ES
\documentclass[12pt]{article}
\usepackage{graphicx}
\usepackage{moodle}
\begin{document}
\begin{quiz}{Matemáticas/Números Complejos}
 
 
 	\begin{multi}[shuffle=true,numbering=none, feedback={	
Si desarrollamos $(1+i)^2=1^2+i^+2i =0+2i$	
	 	}]{Números complejos}
	 	La parte real de $(1+i)^2$ es
		\item* $0$ 
		\item $1$
		\item $2$
		\item $2i$
 	\end{multi}
 
	 
	 \begin{multi}[shuffle=true,numbering=none,feedback={Calcúlese el cuadrado}]%
	 {Números complejos}
		La parte imaginaria de $(1+i)^2$ vale
		\item $2i$
		\item $1$
		\item* $2$
		\item $0$
	 \end{multi}

	 \begin{multi}[shuffle=true,numbering=none]%
	 {Números complejos}
		El módulo de $5+i$ vale 
		\item $6$
		\item* $\sqrt{26}$
		\item $\sqrt{6}$
		\item $\sqrt{24}$
	 \end{multi}

	 \begin{multi}[shuffle=true,numbering=none]%
	 {Números complejos}
		El módule de $i^3$ vale
		\item* $1$
		\item $-1$
		\item $i$
		\item $-i$
	 \end{multi}
	 
	 \begin{multi}[shuffle=true,numbering=none]%
	 {Números complejos}
		El argumento de $1+i$ es
		\item* $\pi/4$
		\item $0$
		\item $\pi/2$
		\item $\sqrt{2}$
	 \end{multi}	 
 			
	 \begin{multi}[shuffle=true,numbering=none]%
	 {Números complejos}
		Resuelve
		\item* $-\frac{1}{2}\pm \frac{\sqrt{3}}{2}\,i $
		\item No tiene solución
		\item $-\frac{1}{2}\pm \frac{3}{2}\, i$
	 \end{multi}
	 
	 \begin{multi}[shuffle=true,numbering=none]%
	 {Números complejos}
		Calcula $\frac{1+i}{1-i}$
		\item $1-i$
		\item* $i$
		\item $1+i$
		\item $-i$
	 \end{multi}	
	 
	 \begin{multi}[shuffle=true,numbering=none]%
	 {Números complejos}
		Calcula el argumento de $9-9\sqrt{3}i$
		\item $45$
		\item $-\pi/4$
		\item $\pi/3$
		\item* $-\pi/3$
	 \end{multi}	  
	 			
\end{quiz}
\end{document}
 
